% 367_Phase_Wahrscheinlichkeit_De_ch.tex
% Johann Pascher, 19. September 2026
% Warum theta = p0 = 2/9: Wahrscheinlichkeit und Phase als zwei Lesarten eines Quotienten

\providecommand{\BM}{\textbf{[B]}}
\providecommand{\KM}{\textbf{[K]}}
\providecommand{\SM}{\textbf{[S]}}
\providecommand{\SetzM}{\textbf{[SETZUNG]}}
\providecommand{\QM}{\textbf{[Q]}}
\providecommand{\EM}{\textbf{[E]}}
\providecommand{\XM}{\textbf{[X]}}

\chapter*{Warum $\theta = p_0 = 2/9$: Wahrscheinlichkeit und Phase als zwei Lesarten eines Quotienten}
\addcontentsline{toc}{chapter}{Warum $\theta = p_0 = 2/9$}

\begin{abstract}
Der Koide-Winkel $\theta=2/9$ (Dok.~292) und die ikosaedrische Projektion
$p_0=|\langle v_0\mid R_5\,v_e\rangle|^2=2/9$ (Dok.~293) entstammen zwei
voneinander unabhängigen Rechnungen und treffen exakt dieselbe rationale Zahl.
Dieses Dokument benennt den Grund und zieht die vollständige Konsequenz.

Der Grund: beide Größen sind derselbe Quotient $2/3^2$ -- Zähler = Anzahl der
nicht-trivialen $\mathbb{Z}_3$-Moden, Nenner = Quadrat der $\mathbb{Z}_3$-Ordnung.
$\theta$ liest ihn als Phase des Massenoperators, $p_0$ als Betragsquadrat
einer algebraischen Projektion. Das ist die Born-Regel-Struktur: dieselbe Zahl,
einmal Betrag und einmal Betragsquadrat desselben Matrixelements.

Die Konsequenz: das einzige Matrixelement $A=\langle v_0\mid R_5\,v_e\rangle$
der $A_5$-Einbettung bestimmt die Koide-Formel vollständig und parameterfrei.
Sein Betrag liefert den Amplitudenkoeffizienten $\sqrt2 = 3|A|$, sein
Betragsquadrat liefert die Phase $\theta = |A|^2 = 2/9$:
\[
  a_k = 1 + 3\,|A|\cdot\cos\!\bigl(|A|^2 + \tfrac{2\pi k}{3}\bigr),
  \qquad A = \langle v_0\mid R_5\,v_e\rangle.
\]
Kein freier Parameter. Kein externes $\theta$. Die Massenamplituden der
geladenen Leptonen folgen direkt aus der $A_5$-Geometrie.

Was transzendent bleibt: $\cos(2/9)$, da $2/9\notin\pi\mathbb{Q}$. Exakte
algebraische Massenamplituden sind daher prinzipiell ausgeschlossen.
\end{abstract}

\section*{Zur Sprache: Warum „Wahrscheinlichkeit"?}
\addcontentsline{toc}{section}{Zur Sprache: Warum „Wahrscheinlichkeit"?}

Der Ausdruck $p_0 = |\langle v_0 \mid R_5\, v_e\rangle|^2$ heißt in der
Quantenmechanik \emph{Übergangswahrscheinlichkeit} -- geborgtes Vokabular,
das irreführen kann. Hier ist weder ein statistisches Experiment noch eine
Häufigkeit gemeint. $p_0$ ist das \emph{Betragsquadrat einer algebraischen
Projektion}: $v_e$ und $v_0$ sind orthonormale Eigenmoden des
$\mathbb{Z}_3$-Zirkulanten, $R_5\in A_5$ ist eine exakte Drehmatrix.
Das Skalarprodukt $\langle v_0\mid R_5\,v_e\rangle$ folgt vollständig
aus der Geometrie von $A_5$ -- kein Zufallselement, keine Häufigkeit.

\textbf{Der eigentliche Befund ist algebraisch:} dieselbe rationale Zahl
$2/3^2$ erscheint in zwei algebraischen Rollen -- als Betragsquadrat einer
Projektion und als Positionsparameter des Massenoperators.
Die Sprache wechselt; das Objekt ist dasselbe.

\paragraph{Externe Parallele \EM.}
Derselbe Quotient $2/3^2$ tritt unabhängig in der konformen Feldtheorie
von $\mathbb{Z}_3$-Orbifolds auf: die beiden Twist-Felder $\sigma_1,\sigma_2$
haben konforme Gewichte $h_k = k(3-k)/(2\cdot3^2) = 1/9$, also
$h_1+h_2 = 2/9$ (Dixon, Friedan, Martinec, Shenker 1987). Auch dort steht
im Zähler die Zahl der nicht-trivialen $\mathbb{Z}_3$-Sektoren und im Nenner
$|\mathbb{Z}_3|^2$. Das ist ein \emph{anderes} mathematisches Objekt --
eine Summe konformer Dimensionen, kein Betragsquadrat einer Projektion --
und es entstammt einer anderen $\mathbb{Z}_3$-Wirkung: dem zentralen
Element $\omega\cdot\mathbb{1}$ auf $\mathbb{C}^3$, nicht der zyklischen
Permutation $C_3\in A_5$ mit Eigenwerten $(1,\omega,\omega^2)$. Die
Parallele belegt lediglich, dass $2/3^2$ eine $\mathbb{Z}_3$-strukturelle
Zählgröße ist und keine FFGFT-spezifische Zahl. Eine Identifikation der
beiden Objekte ist damit nicht behauptet.



\paragraph*{Epistemische Marker:}\mbox{}\\[2pt]
\begin{tabular}{@{}ll@{}}
\textbf{[B]} & algebraisch bewiesen (Python-Prüfskript, exakte Arithmetik) \\
\textbf{[K]} & numerisch bestätigt \\
\textbf{[S]} & strukturell plausibel, Identifikation offen \\
\textbf{[E]} & etablierte externe Mathematik oder Physik \\
\end{tabular}

\medskip
\begin{tabular}{@{}lp{11cm}@{}}
\toprule
\textbf{Symbol} & \textbf{Bedeutung} \\
\midrule
$\mathbb{Z}_3$ & zyklische Gruppe der Ordnung~3; Trialität auf $T^4$ \\
$v_0, v_1, v_2$ & Fourier-Eigenmoden des $\mathbb{Z}_3$-Zirkulanten auf $\mathbb{C}^3$ \\
$\omega$ & $e^{2\pi i/3}$, primitive dritte Einheitswurzel \\
$A_5$ & Ikosaedergruppe; enthält $C_3$ als Untergruppe (Dok.~285) \\
$R_5$ & Fünffachdrehung um $72^\circ$ in $A_5$ \\
$\theta$ & Koide-Phase; $\sqrt{m_k}=M\,(1+\sqrt2\cos(\theta+2\pi k/3))$ \\
$p_0$ & $\lvert\langle v_0\mid R_5 v_e\rangle\rvert^2$, Übergangswahrscheinlichkeit \\
$Q$ & Koide-Invariant $(\sum m_k)/(\sum\sqrt{m_k})^2 = 2/3$ \\
\bottomrule
\end{tabular}

\section{Ausgangslage: zwei Zeugen, ein Wert}

Dok.~293 hält den Befund so fest:

\begin{center}
\begin{tabular}{@{}p{0.30\linewidth}p{0.36\linewidth}p{0.24\linewidth}@{}}
\toprule
\textbf{Zugang} & \textbf{Herkunft} & \textbf{Wert}\\
\midrule
Koide (phänomenologisch) & gemessene Lepton-Massen; keine Geometrie & $\theta=2/9$\\
$C_3$-in-$A_5$ (geometrisch) & $5$-fach-Umverteilung der Moden; keine Massen & $p_0=2/9$\\
\bottomrule
\end{tabular}
\end{center}

Beide Wege wissen nichts voneinander und treffen exakt dieselbe rationale Zahl.
Die Härtetests in Dok.~293 (200 Zufallsachsen, Winkelscan, Ordnungsreihe) zeigen:
$2/9$ ist ikosaeder-spezifisch, kein generischer Wert. Dok.~293 schließt mit dem Satz,
$\theta$ sei ein \enquote{Invariant der Übersetzung} zwischen Differential- und
Spektralgeometrie -- und lässt offen, warum gerade $\theta = p_0$ gilt und nicht etwa
$\theta = \sqrt{p_0}$ oder eine andere Funktion. Diese Lücke wird hier geschlossen. Abschnitt~5 geht einen Schritt weiter: die Koide-Formel wird vollständig in Matrixelementen von $R_5$ geschrieben, sodass die Massenamplituden $a_k$ direkt aus der $A_5$-Geometrie folgen -- ohne freien Parameter und ohne externes $\theta$.

\section{Der gemeinsame Quotient $2/3^2$}

\subsection{$p_0$ als Zählung im $\mathbb{Z}_3$-Modenraum}

Die drei Lepton-Moden sind die Fourier-Eigenmoden des $\mathbb{Z}_3$-Zirkulanten auf der
$\mathbb{C}^3$-Faser (Dok.~282, 288):
\[
  v_0 = \tfrac{1}{\sqrt3}(1,1,1),\qquad
  v_{1,2} = \tfrac{1}{\sqrt3}(1,\omega^{\pm1},\omega^{\pm2}).
\]
Eine Mode ist trivial ($v_0$, symmetrisch unter $\mathbb{Z}_3$), zwei sind nicht-trivial
($v_1,v_2$, galoiskonjugiert). Die Elektronmode ist eine der beiden nicht-trivialen (welche,
entscheidet das Vorzeichen des Winkels, Dok.~293).

Wird $v_e = v_1$ mit $R_5$ gedreht, verteilt sich die Mode auf alle drei Richtungen mit den
Gewichten (Dok.~293, \BM):
\[
  p_0 = \frac{2}{9},\qquad
  p_2 = \frac{5-3\varphi}{9},\qquad
  p_1 = \frac{2+3\varphi}{9},\qquad
  p_0+p_1+p_2 = 1.
\]
Alle drei Nenner sind $9 = 3^2$. Für $p_0$ ist der Zähler die kleine ganze Zahl~$2$ --
die Anzahl der nicht-trivialen Moden. In Worten:
\[
  p_0 = \frac{\#\,\text{nicht-triviale }\mathbb{Z}_3\text{-Moden}}{|\mathbb{Z}_3|^2}
  = \frac{2}{3^2}.
\]
Der Nenner ist das Quadrat, nicht die Gruppenordnung selbst, weil $p_0$ ein
Betragsquadrat einer Amplitude ist: $|\langle v_0|R_5 v_e\rangle| = \sqrt{2}/3$, und
$(\sqrt2/3)^2 = 2/9$. Der Faktor~$3$ im Nenner der Amplitude ist die Normierung der
Moden ($1/\sqrt3$ pro Vektor, zwei Vektoren im Skalarprodukt), der Faktor $\sqrt2$ im Zähler
entstammt der ikosaedrischen Geometrie (Achsenwinkel $37{,}38^\circ$ zwischen $3$-fach und
$5$-fach, Dok.~293).

\subsection{$\theta$ als Phase im $\mathbb{Z}_3$-Modenraum}

Die Koide-Parametrisierung (Dok.~292)
\[
  \sqrt{m_k} = M\,\bigl(1+\sqrt2\cos(\theta + 2\pi k/3)\bigr),\qquad k=0,1,2,
\]
hat zwei Strukturzahlen: die Amplitude $\sqrt2$ und die Phase $\theta$. Dok.~353 zeigt, dass
$\sqrt2$ aus der $\mathbb{Z}_3$-Gleichverteilung folgt \BM. Dok.~291 zeigt, in welchem Kanal
$\theta$ lebt: nicht im symmetrischen Betragsfunktional (dort ist $Q = 2/3$ für jedes $\theta$),
sondern im Phasenkanal der $\mathbb{Z}_3$-Holonomie.

Der empirische Wert ist $\theta = 2/9$ (Dok.~292: Abweichung des PDG-Rückwerts von $2/9$
nur $0{,}9\sigma$, $\tau$-limitiert; die genaue Nachkommastelle hängt vom verwendeten
$m_\tau$-Zentralwert ab, die $\sigma$-Bewertung nicht). Als Bruch gelesen: $\theta = 2/3^2$. Auch hier
steht der Zähler~$2$ für die zwei nicht-trivialen Moden und der Nenner $3^2$ für das Quadrat der
$\mathbb{Z}_3$-Dimension.

Warum das Quadrat? Weil $\theta$ die Phase auf dem \emph{Produktraum} der Moden misst -- nicht
auf der Gruppe $\mathbb{Z}_3$ (dort wäre der natürliche Bruch $2/3$, der Anteil nicht-trivialer
Elemente), sondern auf $\mathbb{C}^3 \otimes \mathbb{C}^3$, dem Raum der Massenoperator-Elemente.
Der hermitesche Zirkulant $M_{ij}$ hat $3^2 = 9$ Einträge, und die Phase, die $M$ von der
Symmetrieachse abweichen lässt, ist ein Bruchteil dieses $9$-dimensionalen Raums.

\subsection{Dieselbe Asymmetrie, zweimal gelesen}

Beide Größen zählen dasselbe:
\begin{equation}\label{eq:quotient}
  \theta = p_0 = \frac{\#\,\text{nicht-triviale }\mathbb{Z}_3\text{-Moden}}{|\mathbb{Z}_3|^2}
  = \frac{2}{3^2} = \frac{2}{9}.
\end{equation}
$p_0$ liest diesen Quotienten als \emph{Wahrscheinlichkeit}: Wie viel der Elektronmode landet
unter der ikosaedrischen Drehung in der symmetrischen Richtung? $\theta$ liest ihn als
\emph{Phase}: Wie weit ist die Massenanordnung von der vollen $\mathbb{Z}_3$-Symmetrie entfernt?
Beide Male wird die Asymmetrie der Lepton-Konfiguration gegenüber der vollen
$\mathbb{Z}_3$-Symmetrie gemessen -- einmal als Anteil, einmal als Winkel.

\section{Die Struktur der Born-Regel}

Das ist die Struktur der Quantenmechanik. Die Born-Regel \EM{} verbindet eine komplexe Amplitude
$A = \langle\psi|\phi\rangle$ mit einer messbaren Wahrscheinlichkeit:
\[
  P = |A|^2.
\]
Die Phase von $A$ ist nicht direkt messbar; nur $|A|^2$ ist es. Dennoch trägt die Phase die
Interferenzinformation -- sie wird sichtbar, sobald zwei Amplituden überlagert werden.

In der Koide-Struktur liegt derselbe Doppelcharakter vor, in einer besonders einfachen Form:
\begin{itemize}
  \item Die Amplitude der ikosaedrischen Projektion ist $|\langle v_0|R_5 v_e\rangle| = \sqrt2/3$
    (Betrag; kein wählbarer Parameter, Dok.~293).
  \item Ihr Quadrat ist die Wahrscheinlichkeit $p_0 = 2/9$.
  \item Der Koide-Winkel $\theta = 2/9$ ist \emph{numerisch derselbe Quotient}, gelesen als
    Phase.
\end{itemize}
Der Zusammenhang ist nicht $\theta = |A|$ (das wäre $\sqrt2/3 \approx 0{,}471$) und nicht
$\theta = \arg A$ (das wäre ein Drehwinkel der Projektion), sondern $\theta = |A|^2 = p_0$.
Das ist der Kern: \emph{die Phase der Massenanordnung ist gleich der Wahrscheinlichkeit der
Modenumverteilung}, weil beide den Quotienten~\eqref{eq:quotient} realisieren.

Dass dieselbe Zahl in beiden Rollen erscheint, ist keine Zufälligkeit der Zahlenwerte, sondern
eine Eigenschaft des $\mathbb{Z}_3$-Modenraums: der Nenner $3^2$ ist in beiden Rollen erzwungen
-- als Normierungsquadrat in $p_0$, als Dimension des Operatorraums in $\theta$. Der Zähler~$2$
ist in beiden Rollen die Anzahl der nicht-trivialen Moden. Das ist die algebraische Gestalt
dessen, was Dok.~293 als \enquote{Invariant der Übersetzung} beschreibt: die Zahl gehört der
Struktur, nicht der Darstellung.

\section{Numerische Bestätigung}

Das Prüfskript \texttt{pruef\_367\_phase\_wahrscheinlichkeit.py} (2/python/Dok367\_Skripte)
bestätigt \KM:
\begin{enumerate}
  \item $p_0 = |\langle v_0|R_5 v_1\rangle|^2 = 2/9$ auf 40 Stellen (mpmath), mit der
    Standardeinbettung $C_3 < A_5$ (Achse $(0,1,\varphi)$).
  \item $|\langle v_0|R_5 v_1\rangle| = \sqrt2/3$ exakt.
  \item $\theta_{\rm emp}$ aus den PDG-Massen (Rückrechnung via $\cos\theta$ aus
    $\sqrt{m_\tau},\sqrt{m_e},\sqrt{m_\mu}$; $m_\tau = 1776{,}86 \pm 0{,}12$~MeV): $0{,}222270$,
    Abweichung von $2/9$: $4{,}8\cdot10^{-5}$; die $\tau$-limitierte Toleranz ist
    $5{,}3\cdot10^{-5}$, die Abweichung entspricht $0{,}9\sigma$ (wie Dok.~292).
  \item $\theta_{\rm emp} = p_0$ innerhalb der $\tau$-Toleranz (empirisch); exakt $= 2/9$
    beiderseits (geometrisch).
  \item Die Alternativen $\theta = \sqrt{p_0}$ ($0{,}471$) und $\theta = \arg\langle v_0|R_5 v_1\rangle$
    treffen den empirischen Winkel nicht.
\end{enumerate}

\section{Die direkte Übersetzung: vom U-Matrixelement zu den Leptonmassen}

Diese Gleichung stellt den vollständigen Schulterschluss her. Sie verbindet
die Geometrie der $A_5$-Einbettung ohne jeden Zwischenschritt oder freien
Parameter direkt mit den Massenamplituden der geladenen Leptonen \BM:
\begin{equation}\label{eq:direkt}
  \boxed{a_k = 1 + 3\,|\langle v_0 \mid R_5\, v_e\rangle|\cdot
    \cos\!\Bigl(|\langle v_0 \mid R_5\, v_e\rangle|^2 + \tfrac{2\pi k}{3}\Bigr).}
\end{equation}

\paragraph{1. Der Amplitudenkoeffizient ($\sqrt2$).}
Der Betrag des Übergangselements ist
$|\langle v_0 \mid R_5\, v_e\rangle| = \sqrt{2/9} = \sqrt2/3$.
Wird dieser Betrag mit der Dimension der $\mathbb{Z}_3$-Symmetrieachse (3)
skaliert, entsteht exakt der Koide-Faktor:
\[
  3\cdot|\langle v_0 \mid R_5\, v_e\rangle| = 3\cdot\frac{\sqrt2}{3} = \sqrt2.
\]

\paragraph{2. Die Phasenverschiebung ($\theta$).}
Nach der Bornschen Regel ist das Betragsquadrat desselben Matrixelements die
Projektionswahrscheinlichkeit~$p_0$:
\[
  p_0 = |\langle v_0 \mid R_5\, v_e\rangle|^2 = \frac{2}{9}.
\]
Indem dieses Quadrat direkt als Phase in den Kosinus eingeht
($\theta = p_0 = 2/9$ rad), tritt die Bornsche Regel als geometrische
Kopplung auf: Amplitude und Phase desselben Matrixelements bestimmen gemeinsam
die Massenstruktur.

\paragraph{3. Die Massenamplituden.}
Für $k\in\{0,1,2\}$ ergeben sich ohne weitere Annahmen die drei Koide-Amplituden \KM:
\begin{align*}
  \text{Tau }(\tau,\;k=0): &\quad a_0 = 1+\sqrt2\cos\!\bigl(\tfrac{2}{9}\bigr)
    \approx 2{,}3794, \\
  \text{Elektron }(e,\;k=1): &\quad a_1 = 1+\sqrt2\cos\!\bigl(\tfrac{2}{9}+\tfrac{2\pi}{3}\bigr)
    \approx 0{,}0404, \\
  \text{Myon }(\mu,\;k=2): &\quad a_2 = 1+\sqrt2\cos\!\bigl(\tfrac{2}{9}+\tfrac{4\pi}{3}\bigr)
    \approx 0{,}5802.
\end{align*}
Daraus folgen unmittelbar die bekannten Massenverhältnisse $m_k \propto a_k^2$.

\paragraph{Einordnung.}
Damit ist die offene Frage beantwortet:
\begin{itemize}
  \item Das Matrixelement $\langle v_0 \mid R_5\, v_e\rangle$ steuert sowohl
    die Amplitude ($\propto |A|$) als auch die Phase ($\propto |A|^2$).
  \item Es braucht kein ungebundenes $\theta$ als externe Eingabe mehr -- die
    Zahl $2/9$ generiert sich im Funktionsaufruf selbst aus dem Betrag des
    Matrixelements.
  \item Die prinzipielle Transzendenz über die Kosinus-Funktion ($\cos(2/9)$
    ist transzendent, da $2/9\notin\pi\mathbb{Q}$) \EM\ bleibt die mathematische
    Schnittstelle, weshalb die Massenverhältnisse irreduzibel transzendente
    Zahlen sind.
\end{itemize}

\section{Was erklärt ist und was offen bleibt}

\paragraph{Erklärt.} Die Übereinstimmung $\theta = p_0$ ist kein Konvergenzzufall zweier
unabhängiger Rechnungen, sondern die zweifache Realisierung desselben Quotienten $2/3^2$ --
einmal als Wahrscheinlichkeit (Betragsquadrat einer normierten Amplitude), einmal als Phase (Anteil
des Operatorraums). Der Grund, warum es $|A|^2$ und nicht $|A|$ ist, liegt im Nenner: der
Modenraum ist $\mathbb{C}^3$, der Operatorraum $\mathbb{C}^3\otimes\mathbb{C}^3$, und die Phase
sitzt im Operatorraum.


\paragraph{Einordnung.} Dass Phase und Wahrscheinlichkeit dieselbe Zahl tragen, ist die Sprache,
in der FFGFT und die Quantenmechanik dieselbe Struktur teilen: Amplitude, Betragsquadrat,
Interferenz. FFGFT gibt dieser Struktur eine geometrische Adresse -- $2/3^2$ auf dem
$\mathbb{Z}_3$-Modenraum -- die Quantenmechanik gibt ihr die Messvorschrift. Beide sagen dasselbe,
in verschiedener Sprache.
